%\begin{marginfigure}[-1cm]
%    \includegraphics[width=5cm]{images/jacques_hadamard.jpg}
%    \caption{Jacques \textsc{Hadamard}}
%\end{marginfigure}

\begin{theo}{Inégalité d'\textsc{Hadamard}}
    Soit $M \in \M_n(\C)$ et soient $X_1, \dots, X_n$ ses vecteurs colonnes. Alors,
    $$|\mathrm{det} M | \leqslant \prod_{i=1}^{n} \Vert X_i \Vert$$
    avec égalité si et seulement si la famille $(X_i)_{1 \leqslant i \leqslant n}$ est orthogonale.
\end{theo}

%\begin{marginfigure}[-2cm]
%    \resizebox{6.5cm}{6.5cm}{%
    \begin{tikzpicture}
        \node[block] (gram) {Déterminant \\ de \textsc{Gram}};
        \node[block, right of=s1] (iwasama) {Décomposition \\ d'\textsc{Iwasama}};
        \node[block, below right of=gram] (hadamard) {Inégalité \\ d'\textsc{Hadamard}};
    
        \draw (gram) edge[bend right, above left] node {...} (hadamard);
    
        \draw (iwasama) edge[bend left] node {$M=OT$} (hadamard);
    
    \end{tikzpicture}    
}
%\end{marginfigure}

Voyons deux démonstrations de ce résultat; une première en utilisant la décomposition d'\textsc{Iwasama} et une deuxième le \refthm{inegalite_gram}.

\begin{preuve}
    \begin{itemize}
        \item Si la matrice $M$ n'est pas inversible alors $\det M = 0$ et comme une norme est à valeurs positives, le résultat est immédiat. 
        \item Supposons que $M$ est inversible. D'après la décomposition d'\textsc{Iwasama}, il existe une matrice $O \in \mathscr{O}_n(\C)$ et $T \defeq (t_{i,j})_{1 \leqslant i, j \leqslant n}$, triangulaire supérieure dont les coefficients diagonaux sont strictement positifs telles que $M = OT$. D'après la multiplicité du déterminant, 
        $$\det M = \det(O) \det(T).$$
        Or $\det O = \pm 1$ donc
        \begin{equation} \label{det}
            |\det M | = |\det T | = \prod_{i=1}^{n} |t_{i,i}|.
        \end{equation}
        Par construction, pour tout $i \in \llbracket 1, n \rrbracket, t_{i,i} \defeq \langle X_i, O_i \rangle$ où $O_i$ est un vecteur unitaire. D'après l'inégalité de \textsc{Cauchy}-\textsc{Schwarz}, pour tout $i \in \llbracket 1, n \rrbracket$, 
        $$|t_{i,i}| = |\langle X_i, O_i \rangle | \leqslant \norme{X_i} \underbrace{\norme{O_i}}_{=1}.$$
        Ainsi d'après (\ref{det}), 
        $$|\det(M)| \leqslant \prod_{i=1}^n \norme{X_i}.$$
    \end{itemize}
        \textcolor{red}{cas d'égalité}
\end{preuve}

\begin{preuve}
    \begin{itemize}
        \item Si la matrice $M$ n'est pas inversible, le résultat est immédiat. 
        \item Supposons que $M$ est inversible. On a 
        $$\Trsp{M} M = \Gram(X_1, \dots, X_n),$$ 
        la matrice de \textsc{Gram} de la famille des colonnes de la matrice $M$. 
        En composant cette relation par le déterminant, 
        $$\det(\Trsp{M}M) = \det \big( \Gram(X_1, \dots, X_n) \big) = \det(M)^2 $$
        car $\det(\Trsp{M}) = \det M$.
        D'après le (\ref{inegalite_gram}), 
        \begin{align*}
            \det \Gram(X_1, \dots, X_n) &\leqslant \prod_{i=1}^n \norme{X_i}^2 \\
            \text{soit } \det (M)^2 &\leqslant \prod_{i=1}^n \norme{X_i}^2
        \end{align*}
        En passant à la racine on obtient l'inégalité d'\textsc{Hadamard}.
    \end{itemize}
\end{preuve}

\marginnote[0cm]{
    \begin{defi}{Parallélotope}
        Soit $(x_1, \dots, x_n)$ une famille libre. Le parallélotope engendré par cette famille est défini par
        $$P \defeq \left\{ x = \sum_{i=1}^n t_i x_i,\ \forall i\ t_i \in [0,1]\right\}.$$
    \end{defi}
}

L'inégalité d'\textsc{Hadamard} nous apprend que le volume du parallélotope défini par les vecteurs colonnes est inférieur ou égal au produit des normes de ses vecteurs et il y a égalité si et seulement si la matrice est diagonale, ou encore que le parallélotope est rectangle. 

\begin{prop}{}
    Soient $\mathscr{S}_n ^{++} (\R)$ l'ensemble des matrices réelles d'ordre $n$ symétriques à valeurs propres strictement positives et $A = (a_{i,j}) \in \mathscr{S}_n ^{++} (\R)$. Alors,
    $$\det(A) \leqslant \prod_{i=1}^{n} a_{i,i}.$$
\end{prop}

\begin{exercice}    
\marginnote[0cm]{Source : exercice 4, TD 14 \cite{acamanes}}
\begin{enumerate}
    \item Soit $(\gamma_1, \dots, \gamma_n) \in (\Re)^n$. Montrer que $B = (\gamma_i \gamma_j a_{i,j}) \in \mathscr{S}_n^{++}(\R)$. 
    \item Montrer que $\det(A)^{1/n} \leqslant \frac{\Tr(A)}{n}$. \\
    \emph{On pourra utiliser l'inégalité arithmético-géométrique}.
    
    \marginnote[-2cm]{
    	\begin{prop}{Inégalité arithmético-géométrique}
            Soient $n \in \Ne$ et $x_1, \dots, x_n$ des réels positifs. Alors, 
            $$\frac{x_1 + \cdots + x_n}{n} \geqslant \sqrt[n]{x_1 \cdots x_n}.$$
            Il y a égalité si et seulement si tous les $x_i$ sont égaux.
        \end{prop}
        Pour d'autres inégalités, lire le chapitre 16, p.117 de la deuxième édition de Raisonnements divins (en fr)
    }
    \item Montrer que pour tout $i \in \llbracket 1, n \rrbracket, a_{i,i} > 0$. On pose $\gamma_i = \frac{1}{\sqrt{a_{i,i}}}$. En déduire l'inégalité d'\textsc{Hadamard}.
\end{enumerate}
\end{exercice}
